% MR25: Original MR24 three-page supplementary methods integrated as one Appendix A.
% All underlying data, equations and counts are preserved.
\section{Technical Methods and Source-Integrity Audit}\label{app:methods}
\noindent\textbf{Status and scope.} This appendix documents deterministic reanalysis of a public annotation file and post-hoc descriptive stability checks. It is \emph{not} independent label validation, a human-coder study, a systematic literature sample, or an estimate of causal severity effects. The original author-annotated rows are not redistributed.

\subsection{Frozen data, provenance, and admission gates}\label{app:provenance}
The unique source is the public annotation object supporting Tang et al., \emph{How Coding Agents Fail Their Users: A Large-Scale Analysis of Developer--Agent Misalignment in 20,574 Real-World Sessions} (EMNLP 2026; DOI \href{https://doi.org/10.18653/2026.emnlp-main.1011}{10.18653/2026.emnlp-main.1011}). Our calculation uses the fixed code-release commit and path:
\par
\noindent\href{https://github.com/ND-SaNDwichLAB/coding-agent-misalignment/blob/c887cc920129b2db0bdf218c8550e98f74f42f7c/data/core/annotated_misalignments.json}{\textbf{Pinned GitHub source file}}.\\
\textbf{Commit:} {\small\ttfamily c887cc920129b2db0bdf218c8550e98f74f42f7c}.\\
\textbf{Path:} \texttt{data/core/annotated\_misalignments.json}.
\par
\noindent\textbf{Source SHA-256:}\\ {\small\ttfamily 7e0693675242a625bdb586bcd8cef23f5da76c65cf840108c3c24b6b479b9f57}.\par The source has 16,118 episode objects. Each composite episode key (source, repository, session, misalignment) is distinct. The observed resolver labels agree structurally with resolution status: the 1,504 records marked RS1 have one RV1/RV2/RV3 resolver; the 14,614 records marked RS2 have null resolver. The source creates its own severity and resolution labels; no reannotation is performed here. Its observational selection rule requires a mismatch made visible by developer correction or pushback, rather than sampling all failures. The source's \href{https://arxiv.org/html/2605.29442v1}{arXiv v1 (28 May 2026)}, Appendix A.2, Table 5, reports on 100 adjudicated records: damage-severity agreement between two source annotators was 0.71, and model-to-gold accuracy was 0.64; the respective resolution-status values were 0.92 and 0.96. These are source-reported per-axis exact-match diagnostics, not a known DS1/DS2 confusion matrix, a full-corpus correction rate, or our own validation. The accepted EMNLP version and released data are distinguished from this archived annotation-quality edition.

\begin{table}[h!]
\centering\small
\caption{Preserved source partitions. RS2 is unknown, not failure.}
\label{tab:appendix-source-partitions}
\begin{tabular}{lrrr}
\toprule
Subset & Episodes & Visible RS1 & Visible RV2\\\midrule
DS1 (effort/trust only) & 14,587 & 1,174 & 1,080\\
DS2 (easily reversible damage) & 1,361 & 330 & 296\\
Other severity classes (DS0/DS3/DS4) & 170 & 0 & 0\\
\midrule
Total & 16,118 & 1,504 & 1,376\\
\bottomrule
\end{tabular}
\end{table}

\noindent\textbf{Observation--classification coupling.} The original codebook requires evidence of an applied, reversible code/state change for DS2, while RS1 requires visible confirmation of correction. Both labels are inferred from the same textual record. This leaves open common-visibility and classification mechanisms even after repository/session matching; the source provides no linked dual-annotator DS1/DS2 confusion matrix to separate them.

\subsection{Units, weights and dependence}\label{app:weights}
Our primary descriptive endpoint is $Y_i=\mathbf{1}\{\mathrm{resolution\_status}_i=\mathrm{RS1}\}$, where $Y=0$ means no recorded visible resolution rather than known failure. In a session $s$ containing both classes, define $D_s=\bar{Y}_{\mathrm{DS2},s}-\bar{Y}_{\mathrm{DS1},s}$. The equal-session estimator is $\bar{D}_{S}=|S|^{-1}\sum_{s\in S} D_s$. The matched sample comprises 853 sessions and 2,775 episodes in 352 source-qualified repositories. The equal-session estimate is $0.1396122927$. The matching gate admits 853/8,590 (9.93\%) sessions and 2,775/15,948 (17.40\%) eligible episodes, across 352/1,229 source-qualified repository keys. The signed difference is positive in 201 sessions, negative in 99, and tied in 553. Of the ties, 528 have zero visible-resolution rates for both severity groups, 22 have rates of one for both, and 3 have equal intermediate rates. The data are nonrandom and this admitted subset is not representative by design.

To preserve clustering beyond the session level, the repository-nested descriptive bootstrap selects 352 source--repository clusters with replacement, keeping \emph{all} matched sessions within each draw. The equal-session mean is recomputed over the resampled sessions (3,000 draws; fixed seed 993720); the 2.5th and 97.5th nearest-order percentiles are $[0.1071711681,0.1729976896]$. The older whole-session resampling interval, $[0.1090339418,0.1702394909]$ (seed 993719), is shown only for comparison because sessions are nested in repositories. The alternative equal-repository estimator gives each of the 352 repository means of session differences equal weight: $0.1405407792$, with repository-resampling percentiles $[0.1017798217,0.1809355907]$.

\begin{table}[h!]
\centering\small
\caption{Nested robustness endpoints (all values in percentage points). Each interval is a descriptive resampling band for this observed nonrandom file.}
\label{tab:appendix-robustness}
\begin{tabularx}{\linewidth}{@{}YrrrY@{}}
\toprule
Restriction/weighting & Sessions & Repos & Mean & Repo-level 2.5--97.5 band\\\midrule
All paired sessions, equal session & 853 & 352 & +13.96 & +10.72 to +17.30\\
All paired sessions, equal repository & 853 & 352 & +14.05 & +10.18 to +18.09\\
Extraction confidence = high & 801 & 339 & +14.30 & +10.87 to +17.66\\
Exactly one episode of each class & 307 & 171 & +16.94 & +11.73 to +22.36\\
\bottomrule
\end{tabularx}
\end{table}

Removing each repository in turn from the 853-session calculation yields an equal-session mean in $[0.1355586404,0.1435967878]$. In the matched sessions, the DS2 rate exceeds the DS1 rate in 201 sessions, is lower in 99, and ties in 553. The original \texttt{confidence} attribute measures support for episode extraction, not certainty of damage-severity coding. Subset comparisons were chosen \emph{after} viewing available fields, are exploratory, and may select a nonrepresentative fraction of the file. No adjustment for an exploratory search over many candidates is claimed; no inferential $p$-values are reported.

\subsection{Logical bounds and assumption frontier}\label{app:bounds}
Let $r_g$ be the visible RV2 share among all episodes in severity group $g$, and $u_g$ the fraction with unknown final status. Define $q_g\in[0,1]$ as a \emph{hypothetical} eventual RV2 fraction among currently unknown cases. Then
\[
\theta_g=r_g+u_gq_g,\qquad
\Delta=\theta_{\mathrm{DS2}}-\theta_{\mathrm{DS1}}=
\left(\frac{296+1031q_2}{1361}\right)-\left(\frac{1080+13413q_1}{14587}\right).
\]
Without assumptions, $\Delta\in[-0.7760687641,0.9009798414]$; the endpoints are logically sharp completions of currently unknown cases, not probabilistic uncertainty limits. The exact zero-difference frontier has the form
\[
q_1=\frac{296/1361-1080/14587}{13413/14587}+\frac{1031/1361}{13413/14587}q_2\approx0.1560+0.8238q_2.
\]

\begin{table}[h!]
\centering\small
\caption{Illustrative \emph{assumption} thresholds, not observed recovery estimates.}
\label{tab:appendix-frontier}
\begin{tabular}{rr}
\toprule $q_2$: unknown DS2 RV2 share & $q_1$ threshold for a tie \\ \midrule
0\% & 15.60\%\\
25\% & 36.20\%\\
50\% & 56.79\%\\
75\% & 77.39\%\\
100\% & 97.98\%\\\bottomrule
\end{tabular}
\end{table}
The special restriction $q_1=q_2=q$ gives $q=0.8856$ as a sign crossing. Without a defensible restriction on $q_1,q_2$, the sign of eventual recovery is not identified. Figure~\ref{fig:frontier} in the main text visualizes this algebra; it is not a fitted empirical curve.

\subsection{Worst-case severity-reclassification certificate}\label{app:relabel}
This check is separate from the missing-eventual-outcome model in Appendix~\ref{app:bounds}. We hold the observed binary RS1/RS2 outcomes fixed and allow the source severity label to move between DS1 and DS2 (no DS0/DS3/DS4 changes). The observed 2-by-2 table is DS1 (RS1=1,174; RS2=13,413) and DS2 (RS1=330; RS2=1,031). To lower the observed DS2-minus-DS1 difference, it can be beneficial to move an RS1 from DS2 to DS1 (count $x$) or an RS2 from DS1 to DS2 (count $y$). Moves in the reverse directions only increase the positive contrast until crossing. Under the allowed changes, the modified comparison is
\[
\frac{330-x}{1361-x+y}\ \leq\ \frac{1174+x}{14587+x-y},\quad 0\leq x\leq330,\quad 0\leq y\leq13413.
\]
We enumerate integer $k=x+y=0,1,\ldots$ and all feasible pairs $0\leq x\leq k$ using exact cross-multiplication, stopping at the first non-positive difference. The first witness occurs at $k=223$, $(x,y)=(223,0)$; all $k<223$ fail the inequality. The witness produces 107/1,138 (9.4025\%) versus 1,397/14,810 (9.4328\%), respectively, and alters 223/15,948 (1.3983\%) original two-class labels; these 223 switches are also 67.58\% of the 330 originally visibly resolved DS2 cases. The latter denominator highlights the concentrated adversarial nature of the witness, not a measured mislabeling prevalence. Restricting the same exercise to each acquisition source separately yields the following exact minima:
\begin{center}
\begin{tabularx}{\linewidth}{@{}lrrrY@{}}
\toprule Stratum & DS1 visible/N & DS2 visible/N & Min. flips & Fraction of eligible labels\\\midrule
Pooled & 1,174/14,587 & 330/1,361 & 223 & 1.3983\%\\
SpecStory & 833/10,861 & 255/1,107 & 170 & 1.4205\%\\
SWE & 341/3,726 & 75/254 & 55 & 1.3819\%\\\bottomrule
\end{tabularx}
\end{center}
The source validation sample does not tell us whether even one such targeted flip is actually warranted. Consequently these minima are \emph{worst-case algebraic perturbation budgets, not estimated annotation error rates, not adjusted results, and not confidence limits}. The observed endpoint depends on classification validity; the unknown eventual recovery endpoint remains separately unidentified even if every severity label were flawless.

\subsection{Reproduction and claim gates}\label{app:reproduction}
Run the scripts with Python 3 (standard library only), after obtaining \path{annotated_misalignments.json} from the pinned author source. Use the exact source file above; the scripts reject any other SHA-256. \emph{From the public GitHub repository root}, use the \texttt{scripts/} paths:
\begingroup\footnotesize
\begin{verbatim}
python scripts/reproduce_visibility_robustness.py \
  --annotations /path/to/annotated_misalignments.json \
  --output visibility_check.json
python scripts/reproduce_nested_robustness.py \
  --annotations /path/to/annotated_misalignments.json \
  --output nested_check.json
python scripts/reproduce_matched_coverage.py \
  --annotations /path/to/annotated_misalignments.json \
  --output coverage_check.json
python scripts/verify_replay.py nested_check.json \
  results/aggregate_results.json
\end{verbatim}
\endgroup
In the \emph{Overleaf source ZIP}, the three analysis scripts sit at the ZIP root rather than in \texttt{scripts/}; omit that prefix for a local extracted-archive replay, and check the original aggregate JSON against the bundled \texttt{MR22\_Aggregate\_Results.json}. The new \texttt{MR26\_Matched\_Coverage\_Results.json} is additional descriptive evidence, not a replacement for the locked original aggregate output.
The JSON output is aggregate-only and includes a deterministic integer-search certificate for minimum severity relabeling budgets. The script enumerates all outcome-concordant DS1/DS2 switches, verifies each candidate by exact integer cross-products, and keeps RS1/RS2 fixed. The complete source records are never copied into the public manuscript or supplementary ZIP. The post-hoc checks are explicitly marked exploratory. Assertions enforce the frozen input, source partitions, 414 mixed-label repositories, 853 mixed-label sessions, 352 repositories supporting those sessions, 801 high-confidence mixed-label sessions, and 307 one-to-one mixed-label sessions.

\noindent\textbf{Interpretive holds.} Neither a repository bootstrap nor matching across recorded sessions can validate the authors' original DS1/DS2 classification, determine RS2 terminal outcomes, or support causal intervention estimates. The original source's 91.49\% post-pushback resolution proportion is valid \emph{only conditional on the visibly resolved 1,504 episodes}. Source-reported validation on the severity axis is materially weaker than on the resolution-status axis in the explicitly identified May 2026 arXiv v1, and no new human reannotation is claimed in this work. Publisher template updates, authorship, CRediT, funding, conflict declaration, exact AI-use disclosure and author-approved persistent public deposit remain separate administrative checks.
